\documentclass[11pt]{article}
\usepackage[margin=.89in]{geometry}
\usepackage{amsmath,amssymb,amsthm}
\usepackage[hidelinks]{hyperref}
\newtheorem{theorem}{Theorem}
\newtheorem{lemma}{Lemma}
\DeclareMathOperator{\diam}{diam}
\title{The sharp inverse-degree bound on mean distance}
\author{Meiji research seat, EPS Round 01}
\date{23 November 2026}
\begin{document}
\maketitle
\begin{abstract}
For fixed $b>2$, the largest mean distance of a connected
$n$-vertex graph with $\sum_v 1/d_v\le b$ is asymptotic to
$(\lceil b\rceil-2)\log n/\log\log n$. The result includes
integer budgets and proves the prediction in Remark 5.2
of the supplied inverse-degree manuscript. The proof
covers the graph by neighborhoods of fewer than $b$
local degree maxima. The worst cost of reaching them is
asymptotic to $\log n/\log\log n$, but their mean access
cost is bounded. We also obtain a second asymptotic term
and an explicit finite upper bound.
\end{abstract}

\section{Statement}
Let $G$ be finite, simple, connected and of order $n\ge2$.
Write $d(v)$ for its vertex degrees,
$J(G)=\sum_v 1/d(v)$, and $\mu(G)$ for the average
shortest-path distance over distinct unordered pairs.
Let $M_b(n)$ be the maximum of $\mu(G)$ over such graphs
with $J(G)\le b$. For each fixed $b>1$ this family is nonempty
at all sufficiently large orders, as the complete graph shows.
Erd\H{o}s, Pach and Spencer \cite{EPS} initiated this extremal
problem in refuting the inequality $\mu(G)\le J(G)$.
Erd\H{o}s--Pach--Spencer use a diameter upper coefficient $6b$
(for fixed $b\ge3$). Dankelmann, Swart and van den Berg \cite{DSB}
improved it to $3b+2$. The supplied manuscript \cite{comp}
gives a matching lower construction for the conjectured mean
coefficient and states the prediction explicitly in Remark~5.2.
The preceding tournament note \cite{baseline} established the
mean coefficient for $2<b\le3$. Neither preceding manuscript
is used as an input to the new upper bound.

\begin{theorem}\label{thm:main}
Fix $b>1$, put $m=\lceil b\rceil-1$ and use natural logarithms.
For $b>2$,
\[
M_b(n)=(m-1)\left(
 \frac{\log n}{\log\log n}
 +\frac{\log n\,\log\log\log n}{(\log\log n)^2}\right)
 +O_b\!\left(\frac{\log n}{(\log\log n)^2}\right).
\]
For $1<b\le2$, $1\le M_b(n)=O_b(1)$ whenever the family
is nonempty. In particular, for every fixed $b>1$,
\[
\lim_{n\to\infty}\frac{M_b(n)}{\log n/\log\log n}
 =\lceil b\rceil-2 .
\]
\end{theorem}

\section{Distance to a local degree maximum}
A vertex $u$ is a \emph{local degree maximum} when
$d(w)\le d(u)$ for every $w$ adjacent to $u$; equality is
allowed. The set $P$ of local degree maxima is nonempty.
Let $\ell(v)=\operatorname{dist}(v,P)$.

\begin{lemma}\label{lem:access}
If $J(G)\le b$, then, with the left side defined to be zero at
$\ell(v)=0$,
\[
\ell(v)\log\frac{\ell(v)}{3b}
 \le\log\frac{n-1}{d(v)}.
\]
Moreover at most $m$ local degree maxima can be pairwise
at distance at least three.
\end{lemma}

\begin{proof}
For $i\ge0$ let $D_i$ be the maximum degree in the closed
radius-$i$ ball centered at $v$. Set $\ell=\ell(v)>0$.
For every $0\le i<\ell$, choose a vertex $w_i$ in the
ball with $d(w_i)=D_i$. It must have distance exactly $i$
from $v$: otherwise all its neighbors would lie in the same
ball and have degree at most $D_i$, making $w_i$ a local
maximum at distance less than $\ell$.
Consequently $N[w_i]$ lies in the three distance layers
$i-1,i,i+1$ (ignoring the negative index).
Each vertex belongs to at most three of the neighborhoods
$N[w_0],\ldots,N[w_{\ell-1}]$.
Every vertex in $N[w_i]$ has degree at most $D_{i+1}$
and $|N[w_i]|=D_i+1$. Summing their inverse-degree charges,
\[
\sum_{i=0}^{\ell-1}\frac{D_i}{D_{i+1}}
 \le\sum_{i=0}^{\ell-1}\sum_{z\in N[w_i]}\frac1{d(z)}
 \le 3b.
\]
The geometric mean of these $\ell$ degree ratios is
$(D_0/D_{\ell})^{1/\ell}\ge(d(v)/(n-1))^{1/\ell}$.
AM--GM proves the claim.

If $u$ is a local maximum,
$\sum_{z\in N[u]}1/d(z)\ge(d(u)+1)/d(u)>1$.
Local maxima pairwise at distance at least three have
disjoint closed neighborhoods. Thus their number is strictly
less than $b$, so it is at most $\lceil b\rceil-1$,
including when $b$ is an integer.
\end{proof}

For $c>0$ and $x\ge0$ let $Q_c(x)$ be the unique $t\ge c$
satisfying $t\log(t/c)=x$. In particular
$Q_c(0)=c$. Put $c=3b$. For $X$ uniform on $V(G)$,
Lemma~\ref{lem:access} and convexity of $t\log(t/c)$
(extended to zero by continuity) give
\[
 \max_v\ell(v)\le Q_c(\log(n-1)),\qquad
 \mathbb E\ell(X)\le Q_c\!\left(\log\frac{b(n-1)}n\right).
\]
For the second bound, Jensen's inequality for $\log$
gives $\mathbb E\log((n-1)/d(X))\le\log((n-1)J(G)/n)$;
Jensen's inequality for $t\log(t/c)$ then gives the assertion,
since whenever $\mathbb E\ell(X)\le c$ the bound already holds.
The argument of $Q_c$ is nonnegative for an admissible graph
because $J(G)\ge n/(n-1)$.

\section{A small weighted spanning tree}
Choose a maximal subset $A\subseteq P$ of local maxima
pairwise at distance at least three. Let $k=|A|\le m$,
and assign each vertex $v$ to a nearest center $a(v)\in A$.
Then $\rho(v)=\operatorname{dist}(v,a(v))\le\ell(v)+2$ by maximality.
Form a graph on $A$, joining two centers if any edge of $G$
joins their classes. It is connected by projecting graph walks.
For such an edge its center-to-center distance in $G$ is
at most $2Q_c(\log(n-1))+5$. Take a spanning tree of
the quotient and put that common length on every tree edge.
The triangle inequality shows that this weighted tree
distance dominates the graph distance between the centers.

Let $X,Y$ be independent uniform vertices and let
$p_e$ be the total probability mass of the centers on one
side of the cut of a tree edge $e$. That edge is crossed
with probability $2p_e(1-p_e)\le1/2$. Thus
\begin{align*}
\mathbb E\,d_G(X,Y)
 &\le 2\,\mathbb E\rho(X)
       +(k-1)\big(Q_c(\log(n-1))+5/2\big).
\end{align*}
Passing to distinct pairs multiplies the left side by
$n/(n-1)$, since the diagonal contributes zero.
In particular every nonempty family at order $n$ satisfies
\begin{equation}\label{eq:finite}
M_b(n)\le\frac n{n-1}\left[
 2Q_{3b}\!\left(\log\frac{b(n-1)}n\right)+4
 +(m-1)\big(Q_{3b}(\log(n-1))+5/2\big)\right].
\end{equation}
For $m=1$ this is $O_b(1)$. The same tree also yields
$\diam(G)\le 2mQ_{3b}(\log(n-1))+5m-1$, without using
the earlier diameter analysis.

For the asymptotic upper bound let $H=\log n$, $K=\log H$,
$U=\log K$. For fixed $c>0$ and a sufficiently large
constant $C_c$, the number $t=H/(K-U-C_c)$ satisfies
$t\log(t/c)>H$ for all large $n$: indeed
$\log(t/c)=K-U-\log c+O_c(U/K)$.
Consequently
\[
 Q_c(\log(n-1))\le H/K+HU/K^2+O_c(H/K^2),
\]
by expanding the denominator, using $U^2/K\to0$.
The other $Q_c$ term in \eqref{eq:finite} is $O_b(1)$.

\section{Matching construction at every large order}
Suppose $b>2$ and $m=\lceil b\rceil-1\ge2$.
Let $A_b>\max(2,8m/(b-m))$ be a fixed constant,
choose a fixed $C_b\ge\log(A_b+1)+2$, and put
\[
r=\left\lfloor\frac{H}{K-U+C_b}\right\rfloor,
\qquad q=\lceil A_b r\rceil .
\]
For all large $n$,
$\log q\le \log r+\log(A_b+1)\le K-U+C_b-1$.
Hence $r\log q\le H-H/(K-U+C_b)\le H-2K$,
so $q^r\le n/H^2$.
Make a sequence of $m$ peaks of size $q^r$, with the gap
between successive peaks descending geometrically to one
and ascending again:
\[
 q^r,q^{r-1},\ldots,q,1,q,\ldots,q^{r-1},q^r .
\]
Consecutive gaps share their terminal peak.
Replace each layer by a clique, add all edges between
consecutive layers, and no other edges. Its initial order is
$O_b(q^r)=O_b(n/H^2)$. Add the remaining vertices to just
the two end layers so that their resulting orders differ by
at most one.

Every peak of size $s$ contributes at most $s/(s-1)$ to
$J(G)$, and every nonpeak contributes at most $1/q$ by
its adjacent larger layer. Thus the total contribution is
at most
\[
 m+O_b(q^{-r})+2(m-1)r/q
 \le m+O_b(q^{-r})+2(m-1)/A_b<b
\]
for all sufficiently large $n$. The two extreme layers each
have order $n/2+O_b(n/H^2)$ and have mutual distance
$2r(m-1)$, since an edge changes layer index by at most one.
Pairs between them give
$\mu(G)\ge(m-1)r(1-O_b(H^{-2}))$.
Finally $r=H/K+HU/K^2+O_b(H/K^2)$, by the same expansion.
For $1<b\le2$, the lower bound 1 is automatic for a
nonempty connected family. This proves Theorem~\ref{thm:main}.

\paragraph{Scope.}
The upper bound is valid uniformly for all graphs at each
fixed admissible $b$. The error constants are allowed to depend
on $b$; no uniform transition as $b$ approaches an integer
with $n$ is asserted. The conclusion is an asymptotic
extremal formula, not exact finite-order extremizer
classification. The elementary argument does not require
computational validation.

\begin{thebibliography}{9}
\bibitem{EPS}
P. Erd\H{o}s, J. Pach and J. Spencer,
\emph{On the mean distance between points of a graph},
Congr. Numer. 64 (1988), 121--124.
\href{https://www.renyi.hu/~p_erdos/1988-23.pdf}{Author archive}.
\bibitem{DSB}
P. Dankelmann, H. C. Swart and P. van den Berg,
\emph{Diameter and inverse degree},
Discrete Math. 308 (2008), 670--673.
\href{https://doi.org/10.1016/j.disc.2007.07.053}{DOI}.
\bibitem{comp}
\emph{The diameter of graphs with given inverse degree:
the constant in a theorem of Erd\H{o}s, Pach and Spencer},
supplied unrefereed manuscript, dated 24 November 2026,
Theorem 1.2 and Remark 5.2.
\bibitem{baseline}
Shaitetsu, \emph{The first two regimes of the inverse-degree
mean-distance problem}, tournament companion note, 23 November 2026.
\end{thebibliography}
\end{document}
